\section{Introduction}

Despite theoretical expectations that denser reward signals should accelerate reinforcement learning, we find that at proof-of-concept scale, all feedback granularities---binary, categorical, and high-bandwidth---produce identical outcomes: zero learning. This finding highlights the importance of adequate experimental scale for code LLM reward ablation studies.

Reinforcement learning from execution feedback (RLEF) has emerged as a promising paradigm for improving code LLMs. Prior work has explored various feedback granularities: binary pass/fail rewards \citep{le2022coderl}, continuous test pass rates \citep{liu2023rltf}, and multi-signal combinations incorporating error type information \citep{shojaee2023ppocoder}. Information theory suggests that higher bandwidth signals---those providing more bits per gradient update---should enable more precise credit assignment and faster convergence.

\citet{liu2023rltf} compared binary and fine-grained feedback in RLTF, finding improvements with richer signals. Our work aimed to provide additional controlled comparison using the ``information bandwidth'' framework. However, our PoC-scale experiment was insufficient to produce any learning signal.

\textbf{Important context:} We use REINFORCE with advantage baseline, not PPO. Prior successful work (CodeRL, PPOCoder) used different algorithms and trained for multiple epochs. Our failure may reflect algorithm choice, scale, or both.

The result was that at proof-of-concept scale (1 epoch, 3 seeds per condition), all nine training runs produced 0\% pass@1 on the evaluation set. No condition learned. The statistical comparison we planned was impossible---there was no variance to analyze.

\textbf{Zero-shot baseline:} CodeLlama-7B-Instruct achieves approximately 30--40\% pass@1 on MBPP in zero-shot evaluation \citep{roziere2023codellama}. Our 0\% post-training result indicates either catastrophic forgetting during training or evaluation methodology differences. This is a critical limitation we did not fully diagnose.

This null result is not a falsification of the information bandwidth hypothesis. It is a finding about our specific experimental setup. Our PoC compute budget ($\sim$94 gradient steps per epoch with batch size 4) was below the threshold at which policy improvement emerges for this configuration.

We report this negative result because it documents:

\begin{enumerate}
    \item \textbf{Setup-specific lower bound.} This specific configuration (REINFORCE, 1 epoch, CodeLlama-7B, MBPP) does not produce learning. We cannot generalize to all code LLM reward ablation without additional evidence.
    \item \textbf{Methodology for replication.} Future researchers can attempt adequate-scale versions of this experiment with the information bandwidth framework.
    \item \textbf{Diagnostic gap.} We lacked training curves, loss plots, and intermediate metrics that would distinguish ``no learning'' from ``catastrophic forgetting'' from ``implementation error.''
\end{enumerate}
